\documentclass[12pt,a4paper]{article}
\usepackage[utf8]{inputenc}
\usepackage[spanish,es-noshorthands]{babel}
\usepackage{amsmath}
\usepackage{siunitx}
\usepackage{circuitikz}
\usepackage[hidelinks]{hyperref}
\usepackage{geometry}
\geometry{a4paper, margin=2.5cm}
\usepackage{microtype}
\emergencystretch 5em
\sloppy

\title{Manual de Arquitectura y Estructura del Proyecto}
\author{Agente IA (Orquestador)}
\date{\today}

\begin{document}
\maketitle
\tableofcontents
\newpage

\section{Introducción}
Este manual describe la estructura de directorios y archivos presentes en el espacio de trabajo de Electrónica. El entorno está diseñado bajo una arquitectura de tres capas, que separa la lógica probabilística (LLMs) de la ejecución determinista (scripts), y también contiene recursos para IoT, Diseño de Circuitos Integrados (EDA) y Redacción Académica.

\section{Directorios Principales}
\begin{itemize}
    \item \textbf{\texttt{.agent/}}: Contiene las instrucciones del sistema para el agente de IA (4 archivos: \texttt{AGENT\_FRAMEWORK.md}, \texttt{AGENT\_INSTRUCTIONS.md}, \texttt{latex.md}, \texttt{python.md}), definiendo su marco operativo, restricciones de entorno y protocolos de actuación.
    \item \textbf{\texttt{directives/}}: Capa 1 de la arquitectura. Contiene manuales de operación (SOPs) en formato YAML (17 archivos). Definen \textit{qué} debe hacerse (ej. \texttt{scrape\_website.yaml}, \texttt{analizar\_imagen.yaml}, \texttt{git\_update.yaml}, \texttt{evaluar\_examen\_estudiante.yaml}, \texttt{evaluar\_practica\_laboratorio.yaml}, \texttt{elaborar\_examen.yaml}, \texttt{elaborar\_ejercicios.yaml}, \texttt{imagen\_a\_kicad.yaml}).
    \item \textbf{\texttt{execution/}}: Capa 3 de la arquitectura. Contiene scripts deterministas en Python con una sola responsabilidad (14 archivos: \texttt{env\_diagnostic.py}, \texttt{scrape\_single\_site.py}, \texttt{analizar\_imagen.py}, \texttt{evaluar\_examen.py}, \texttt{generar\_informe.py}, \texttt{generar\_informe\_imagen.py}, \texttt{alert\_user.py}, \texttt{elaborar\_examen.py}, \texttt{elaborar\_ejercicios.py}, \texttt{generar\_examen\_latex.py}, \texttt{generar\_ejercicios\_latex.py}, \texttt{generar\_kicad\_llm.py}, \texttt{generar\_kicad\_sch.py}, \texttt{extraer\_netlist\_imagen.py}).
    \item \textbf{\texttt{.tmp/}}: Directorio temporal utilizado por la capa de orquestación para almacenar estados de ejecución (como \texttt{run\_state.json}) y artefactos temporales de procesamiento.
    \item \textbf{\texttt{docs/}}: Documentación técnica y académica (21 subdirectorios: CIRC\_DISP\_ELECT/, SMPS/, EDA/, PROTECTOR\_120VAC/, ZBAR\_PRACTICAS/, EASYEDA/, IMAGENES/, RAG/, vLLM/, OPENCODE/, MEDIDOR\_ENERGIA/, PC\_ASUS/, PC\_PARA\_IA/, SERVIDOR\_POWEREDGE\_R610/, SISTEMA\_INTERNAC/, CIRC\_PARA\_RESP\_RAPIDAS/, PROYECTO\_MECATRONICA/, Ventilador\_3\_Velocidades/, KICAD/, AGENTE\_IA/, MANUAL/, entre otros).
    \item \textbf{\texttt{cursos/}}: Material didáctico, notas de cursos de Electrónica, apuntes de Laboratorio y Tesis (6 subdirectorios: DISP\_ELECTRONICOS/, INT\_ELECTRONICA/, TESIS/, PLAN\_ESTUDIOS/, LABORATORIO\_I\_FISICA/, LABORATORIO\_II\_FISICA/). Contiene además subcarpetas \texttt{examen\_estudiante/} e \texttt{informe\_examen/} para el flujo de evaluación.
    \item \textbf{\texttt{directives/rubricas/}}: Rúbricas en formato YAML para la evaluación de prácticas de laboratorio, inyectadas como contexto en el prompt de Gemini.
    \item \textbf{\texttt{chroma\_db/}}: Base de datos vectorial autogenerada, utilizada por el sistema RAG (Retrieval-Augmented Generation) para la memoria y búsqueda semántica del agente.
    \item \textbf{\texttt{Agente\_EDA/}}: Recursos específicos para el agente EDA, como esquemas (schemas) de validación y pruebas para la generación de netlists/BOM hacia EasyEDA.
\end{itemize}

\section{Archivos de Configuración y Scripts Raíz}
\begin{itemize}
    \item \textbf{\texttt{AGENTS.md}}: Reglas globales del proyecto, convenciones de código y estructura general del espacio de trabajo.
    \item \textbf{\texttt{.env}}: Plantilla de variables de entorno con credenciales comentadas (\texttt{GROQ\_API\_KEY}, \texttt{GITHUB\_TOKEN}) como referencia para la configuración local.
    \item \textbf{\texttt{.groq\_api\_key}}: Archivo oculto con la clave API de Groq (excluido de git vía \texttt{.gitignore}), usado por \texttt{rag\_system.py} y \texttt{agent\_eda.py}.
    \item \textbf{\texttt{mcp\_latex\_server.py}}: Servidor MCP local (Layer 2) que expone la herramienta \texttt{compilar\_latex} para automatizar la compilación de reportes académicos.
    \item \textbf{\texttt{mcp\_evaluar\_server.py}}: Servidor MCP local (Layer 2) que expone la herramienta \texttt{evaluar\_examen\_estudiante} con compilación integrada de informes de evaluación a PDF.
    \item \textbf{\texttt{mcp\_elaborar\_server.py}}: Servidor MCP local (Layer 2) que expone la herramienta \texttt{elaborar\_nuevo\_examen} con generación de preguntas y compilación secuencial de examen y solucionario a PDF.
    \item \textbf{\texttt{mcp\_analizar\_server.py}}: Servidor MCP local (Layer 2) que expone la herramienta \texttt{analizar\_imagenes\_circuito} para análisis de visión y compilación de informes técnicos.
    \item \textbf{\texttt{mcp\_diagnostico\_server.py}}: Servidor MCP local (Layer 2) que expone la herramienta \texttt{generar\_diagnostico\_sistema} para telemetría de hardware/software y generación de reportes en PDF.
    \item \textbf{\texttt{.gitignore}}: Exclusiones de seguimiento git: \texttt{chroma\_db/}, entornos virtuales, \texttt{\_\_pycache\_\_}, archivos auxiliares LaTeX, clave API de Groq.
    \item \textbf{\texttt{requirements.txt}}: Dependencias para los scripts deterministas de ejecución: \texttt{psutil} y \texttt{PyYAML}.
    \item \textbf{\texttt{rag\_system.py}}: Script del Chatbot RAG. Vectoriza documentos (TEX, MD, PDF) en ChromaDB y gestiona las respuestas usando el LLM vía Groq.
    \item \textbf{\texttt{agent\_eda.py}}: Agente especializado en EDA. Extrae netlists y BOM de archivos LaTeX (circuitikz) para exportar hacia EasyEDA Standard.
    \item \textbf{\texttt{execution/evaluar\_examen.py}}: Evalúa exámenes y prácticas de laboratorio usando LLM multimodal (Gemini/OpenRouter). Renderiza el PDF en imágenes y envía en bloque al modelo junto con la System Instruction.
    \item \textbf{\texttt{execution/elaborar\_examen.py}}: Genera exámenes escritos completos mediante LLM multimodal, con preguntas, opciones múltiples y solución paso a paso.
    \item \textbf{\texttt{execution/elaborar\_ejercicios.py}}: Genera ejercicios de electrónica con LLM, estructurados en JSON para su posterior formateo LaTeX.
    \item \textbf{\texttt{execution/generar\_examen\_latex.py}} / \texttt{generar\_ejercicios\_latex.py}: Convierten los JSON de examen/ejercicios en documentos LaTeX compilables.
    \item \textbf{\texttt{execution/generar\_informe.py}}: Genera un informe académico profesional en LaTeX a partir del JSON de evaluación producido por \texttt{evaluar\_examen.py}.
    \item \textbf{\texttt{execution/analizar\_imagen.py}}: Analiza una o más imágenes (JPG/PNG) con LLM multimodal (Gemini/OpenRouter) y devuelve un análisis estructurado en JSON.
    \item \textbf{\texttt{execution/generar\_informe\_imagen.py}}: Genera un informe profesional en LaTeX a partir del JSON de análisis producido por \texttt{analizar\_imagen.py}.
    \item \textbf{\texttt{execution/generar\_kicad\_llm.py}} / \texttt{generar\_kicad\_sch.py} / \texttt{extraer\_netlist\_imagen.py}: Flujo EDA para convertir imágenes de circuitos o descripciones a esquemas KiCad.
    \item \textbf{\texttt{flujo\_evaluar\_examen.py}}: Orquestador (Layer 2) que ejecuta el flujo completo de evaluación: analiza el PDF con el LLM, genera el informe LaTeX y notifica al usuario.
    \item \textbf{\texttt{flujo\_elaborar\_examen.py}}: Orquestador (Layer 2) para el flujo de elaboración de exámenes (generar contenido $\to$ convertir a LaTeX $\to$ notificar).
    \item \textbf{\texttt{flujo\_elaborar\_ejercicios.py}}: Orquestador (Layer 2) para el flujo de elaboración de ejercicios de electrónica.
    \item \textbf{\texttt{flujo\_analizar\_imagen.py}}: Orquestador (Layer 2) del flujo de análisis de imágenes: analiza con LLM, genera informe LaTeX, presenta resumen y notifica al usuario.
    \item \textbf{\texttt{flujo\_imagen\_a\_kicad.py}}: Orquestador (Layer 2) que convierte imágenes de circuitos a esquemas KiCad usando LLM.
    \item \textbf{\texttt{flujo\_diagnostico.py}}: Orquestador (Layer 2) del flujo de diagnóstico de sistema: recopila telemetría de hardware y software, genera plantilla LaTeX y compila el informe a PDF.
    \item \textbf{\texttt{clean\_latex.py}}: Utilidad para eliminar archivos auxiliares y temporales de compilación de LaTeX (.aux, .log, .out, etc.).
    \item \textbf{\texttt{fix\_latex.py}}: Script de corrección automática de errores comunes en archivos LaTeX generados por LLM.
    \item \textbf{\texttt{md\_to\_pdf.py} y \texttt{merge\_pdfs.py}}: Utilidades para la conversión y manipulación de documentos PDF.
    \item \textbf{\texttt{ren\_archivos.py}}: Script para el renombrado masivo de archivos eliminando cadenas específicas.
    \item \textbf{\texttt{test\_generator.py}} y \texttt{test\_regex.py}: Pruebas unitarias (sin dependencias externas) para el generador JSON de EasyEDA y expresiones regulares.
    \item \textbf{\texttt{db\_state.json}}: Archivo que lleva el registro del estado y actualizaciones incrementales de la base de datos RAG vectorial.
    \item \textbf{\texttt{opencode.json}}: Archivo de configuración que enlaza las instrucciones del agente alojadas en \texttt{.agent/} (\texttt{\{"instructions": [".agent/*.md"]\}}).
    \item \textbf{\texttt{README.md}}: Documento de presentación del repositorio con instrucciones de instalación, uso y descripción general del proyecto.
    \item \textbf{\texttt{docs/MANUAL/manual.tex}}: El presente manual en formato LaTeX, que describe la arquitectura, estructura y componentes del espacio de trabajo.
    \item \textbf{\texttt{git-update.sh} y \texttt{update\_repo.sh}}: Scripts de automatización en bash para sincronización y actualización del repositorio Git.
\end{itemize}

\section{Servidores MCP Locales (Model Context Protocol)}
El Model Context Protocol (MCP) es un estándar abierto que permite conectar modelos de inteligencia artificial en la nube con recursos locales del sistema de forma segura. En este proyecto se han implementado cinco servidores MCP locales desarrollados en Python mediante el framework \texttt{fastmcp}:

\begin{itemize}
    \item \textbf{Servidor LaTeX (\texttt{mcp\_latex\_server.py})}: Expone la herramienta \texttt{compilar\_latex} para compilar cualquier archivo LaTeX a PDF con resolución de referencias y limpieza en una sola llamada de herramienta.
    \item \textbf{Servidor Evaluador (\texttt{mcp\_evaluar\_server.py})}: Expone \texttt{evaluar\_examen\_estudiante}. Automatiza el análisis de exámenes, generación de informes y compilación a PDF en caliente.
    \item \textbf{Servidor Elaborador (\texttt{mcp\_elaborar\_server.py})}: Expone \texttt{elaborar\_nuevo\_examen}. Genera 5 preguntas con LLM, crea los fuentes LaTeX y compila automáticamente examen y solucionario de forma secuencial.
    \item \textbf{Servidor de Visión (\texttt{mcp\_analizar\_server.py})}: Expone \texttt{analizar\_imagenes\_circuito}. Analiza diagramas o fotos, genera informe LaTeX y compila el documento técnico a PDF.
    \item \textbf{Servidor de Diagnóstico (\texttt{mcp\_diagnostico\_server.py})}: Expone \texttt{generar\_diagnostico\_sistema}. Realiza la telemetría (CPU, RAM, GPU, disco, OS), genera el informe LaTeX y compila el PDF de diagnóstico.
\end{itemize}

La integración de estos servidores proporciona un ahorro significativo de tokens al evitar que el cliente de IA maneje grandes volúmenes de logs de compilación o código fuente en su contexto. Asimismo, permite la automatización total de flujos de trabajo sin intervención manual.

\section{Conclusión}
Esta estructura garantiza una alta modularidad. El agente IA actúa como "middleware" o capa de orquestación (Capa 2), leyendo intenciones, buscando las instrucciones en \texttt{directives/}, delegando el trabajo duro a \texttt{execution/} y documentando hallazgos en \texttt{chroma\_db/} o en archivos \texttt{.tex}. Todo ello cumpliendo con las reglas de codificación y formato estipuladas.

\end{document}
